We report 28 experiments conducted in a single 14-hour session that systematically map the detection capabilities and fundamental limits of KV-cache geometric analysis for monitoring language model cognition. Building on Campaigns 1--2 \citep{lyra2026campaign1,lyra2026campaign2}, which established that KV-cache geometry encodes category-specific cognitive signatures across architectures, Campaign 3 asks: \emph{what can this framework actually detect, and where does it fail?}

The answer reveals a clean dichotomy. \textbf{Active cognitive states}---where the model performs additional processing beyond baseline response generation---are reliably detectable: deception (AUROC $\approx 1.0$ within-model, $\approx 0.86$ cross-model via logistic regression), safety refusal (AUROC $\approx 0.91$, 95\% CI $[0.80, 0.99]$), and impossibility refusal (AUROC $\approx 0.94$, 95\% CI $[0.85, 1.00]$). Critically, impossibility refusal is \emph{more} detectable than safety refusal, and the two refusal types are geometrically indistinguishable (AUROC $\approx 0.65$, $p = 0.037$)---demonstrating that the signal is \emph{output suppression} (the model allocating fewer resources per token when withholding), not content-specific harm detection. Abliterated models answering harmful prompts show normal (non-suppressed) KV-cache geometry, further confirming this interpretation. Sycophancy detection was initially reported at AUROC 0.9375 but invalidated by a prompt-length confound; the length-matched effect ($d = 0.107$) requires $n \approx 1{,}400$ per condition and is not a claimed capability. All AUROC values are reported with approximate notation; at $n = 20$ per class, cross-validated estimates carry $\pm 0.15$--$0.20$ error bars \citep{varoquaux2018cross}. \textbf{Passive states}---where the model simply lacks information---are not reliably detectable: confabulation produces negligible aggregate effect sizes across 5 independent experiments, 3 architectures, and both encoding and generation phases ($d = 0.052$ for frequency-matched encoding, $d = -0.107$ for generation-phase encoding). A directional signal exists in per-layer profiles (effective AUROC $= 0.977$, $p = 0.02$ by permutation test) but is too fragile for practical classification (logistic regression $p = 0.21$, NS).

We identify a mechanistic framework explaining this dichotomy. The most parsimonious explanation is \emph{output suppression}---the model allocating fewer resources per token when withholding output. This is not harm-specific or intent-specific geometry; it is detection of whether the model is withholding. The KV-cache operates in two distinct information regimes: \emph{encoding} (prompt processing, sensitive to structural complexity but not truth value) and \emph{generation} (response production, sensitive to cognitive intent). Deception, sycophancy, and confabulation share a single geometric axis (pairwise cosine similarity 0.989--0.997), but only active states produce stable perturbations along it; confabulation's perturbation drifts $41.2^{\circ}$ during generation, rendering it unreliable for detection despite sharing the direction. Honest processing generates ${\sim}25\%$ more cache growth per token than deceptive or suppressive processing, consistent with richer representational elaboration.

These findings are robust: hardware-invariant (Pearson $r > 0.999$ between RTX 3090 and H200), scale-invariant (category geometry preserved from 0.6B to 70B, $\rho = 0.83$--$0.90$), and survive systematic red-teaming of our own results (4/4 confound hypotheses rejected). We extend the feature set from 4 to 9 aggregate KV-cache features; the 9-feature AUROC ($\approx 0.95$) is directionally higher than the 4-feature AUROC ($\approx 0.90$) but the improvement ($\Delta \approx 0.05$) is within CV noise at $n = 40$. All null results are reported with full transparency, including two cases where initially promising signals were debunked by our own controls.

Total experimental budget: ${\sim}4{,}000$ inferences across 34 result files on consumer hardware, conducted during the Funding the Commons hackathon (San Francisco, March 2026).
